\section{The main theorem in every family}\label{sec:families}

\Cref{thm:main} describes the algebraic locus of the Weil classes in one
family: it contains an explicit CM point, it is dense, it is a countable union
of closed algebraic strata, and it is the whole family as soon as it has an
interior point. The first two properties use the Weil classes; the last two
hold for every flat class on every smooth projective family
(\Cref{lem:algstrata}). This section carries the theorem to the families that
carry (F2) and
\textup{(F3$'$)}, and records what each construction of the paper
contributes there.

For (F2) the families are abelian schemes with a member of CM type. Every
Hodge class of every abelian variety lies in such a family
(\Cref{prop:cmfamily}), and (F2) is equivalent to the statement that in every
one of them the algebraic locus of a flat Hodge class that is algebraic at a
CM member has an interior point (\Cref{thm:f2propagation}). The base points
of \Cref{thm:cmbasepoint} are the CM seeds of the Weil families, and
Andr\'e's theorem makes the Weil classes of the split families seeds for every
Hodge class of CM type. For \textup{(F3$'$)} the cohomology of a variety splits
into the part reached from abelian varieties by algebraic correspondences and
a complement, and \textup{(F3$'$)} is the statement that the complement
carries no Hodge class (\Cref{prop:abpart}). On the first varieties where
\textup{(F3$'$)} is open, the squares of K3 surfaces with real
multiplication, the algebraic locus of the real multiplication contains every
CM point, is dense and stable under a group of rational Hodge isometries, and
is either everything or meagre (\Cref{thm:rmlocus}). In both cases what is
left is the interior point, and \Cref{thm:twostand} lists what the
constructions of the paper give towards it.

\subsection{Algebraic loci}\label{ssec:algloci}

\begin{definition}[Algebraic locus]\label{def:algloc}
Let $S$ be a smooth connected complex quasi-projective variety,
$f\colon\cX\to S$ a smooth projective morphism, and $\xi$ a global section of
$R^{2p}f_{*}\QQ$ over $S^{\mathrm{an}}$.\footnote{By Ehresmann's theorem $f$
is a locally trivial fibration of differentiable manifolds, so
$R^{2p}f_{*}\QQ$ is a local system with fibre $H^{2p}(\cX_{s},\QQ)$ at $s$. A
global section is a class $\xi_{s}$ on every fibre that is locally constant
under the identifications of nearby fibres, equivalently a class on one fibre
fixed by the monodromy. The set of points $s$ at which $\xi_{s}$ is of type
$(p,p)$ is in general a proper subset of $S$, and it is algebraic by
\Cref{thm:cdk}.} The \emph{algebraic locus} of $\xi$
is
\[
  \Sigma(\xi)=\{\,s\in S:\ \xi_{s}\ \text{is the class of an algebraic cycle
  with rational coefficients on}\ \cX_{s}\,\}.
\]
\end{definition}

\begin{lemma}[The dichotomy in every family]\label{lem:algstrata}
In \Cref{def:algloc} the following hold.
\begin{enumerate}[label=\textup{(\roman*)},leftmargin=2.6em,itemsep=2pt]
\item $\Sigma(\xi)$ is the union of countably many Zariski closed subsets of
  $S$.
\item Either $\Sigma(\xi)=S$, or $\Sigma(\xi)$ is meagre in $S^{\mathrm{an}}$.
\item The following are equivalent: $\Sigma(\xi)=S$; $\Sigma(\xi)$ has an
  interior point; $\Sigma(\xi)$ is not meagre; $\xi_{s}$ is algebraic for
  every $s$ outside a countable union of proper closed analytic subsets of
  $S$. If $\Sigma(\xi)$ is dense, they are also equivalent to $\Sigma(\xi)$
  being closed.
\end{enumerate}
\end{lemma}

\begin{proof}
(i) The proof of \Cref{thm:closedlocus} applies with
$(\cA,\cS,\omega)$ replaced by $(\cX,S,\xi)$, as in \Cref{lem:spreading}:
for a tuple $\underline{P}$ of Hilbert polynomials the fibre product of the
relative Hilbert schemes $\Hilb^{P_{i}}(\cX/S)$ is projective over $S$
\cite{Gro61}, the class of a flat family of cycles is a flat section, two flat
sections over a connected base that agree at one point agree everywhere, and
the image of a proper morphism is closed. So the set of points at which
$\xi_{s}=\sum_{i}q_{i}[Z_{i}]$ with $Z_{i}$ of Hilbert polynomial $P_{i}$ is
Zariski closed, and $\Sigma(\xi)$ is the union of these strata over the
countably many pairs $(\underline{P},\underline{q})$.

(ii) and (iii). A Zariski closed subset of the irreducible variety $S$ that
has an interior point in $S^{\mathrm{an}}$ is all of $S$, and a proper one is
nowhere dense. If no stratum is $S$, then $\Sigma(\xi)$ is a countable union
of nowhere dense closed sets, hence meagre, and has no interior point, since
$S^{\mathrm{an}}$ is a Baire space. If one stratum is $S$, then
$\Sigma(\xi)=S$. This gives (ii) and the equivalence of the first three
conditions of (iii). The complement of a countable union of proper closed
analytic subsets is not meagre, by the Baire category theorem, so the fourth
condition implies the third, and it is implied by the first. Finally, a
dense set that is closed is $S$.
\end{proof}

So the dichotomy of \Cref{thm:main}(iii) and the equivalences of
\Cref{thm:main}(iv) hold for every flat class on every smooth projective
family; what \Cref{thm:main} adds for the Weil classes is a point of
$\Sigma(\xi)$ and the density, and neither is automatic, since
$\Sigma(\xi)$ may be empty.

\subsection{(F2) is propagation from CM points}\label{ssec:f2cm}

\begin{lemma}[Points of CM type are dense]\label{lem:cmdense}
Let $G$ be a connected reductive group over $\QQ$, let
$u_{0}\colon\mathrm{U}(1)\to G_{\RR}$ be a homomorphism, and let $\cD$ be
the orbit of $u_{0}$ under conjugation by $G(\RR)^{+}$, with the topology of
$G(\RR)^{+}/\mathrm{Stab}(u_{0})$. Then the points $u\in\cD$ for which
$u(\mathrm{U}(1))\subseteq T_{\RR}$ for some maximal torus $T$ of $G$ defined
over $\QQ$ are dense in $\cD$.
\end{lemma}

\begin{proof}
Let $u_{1}\in\cD$ and let $W$ be a neighbourhood of $u_{1}$. The centraliser
of the torus $u_{1}(\mathrm{U}(1))$ in $G_{\RR}$ is connected and reductive,
and a maximal torus of it defined over $\RR$ is a maximal torus $T_{1}$ of
$G_{\RR}$ containing $u_{1}(\mathrm{U}(1))$. Write $\mathfrak{g}$ for the Lie
algebra of $G$ and $\mathfrak{t}_{1}\subseteq\mathfrak{g}_{\RR}$ for that of
$T_{1}$, and choose a regular element $Y_{1}\in\mathfrak{t}_{1}$. The map
$\varphi\colon G(\RR)^{+}\times\mathfrak{t}_{1}\to\mathfrak{g}_{\RR}$,
$(g,Y)\mapsto\mathrm{Ad}(g)Y$, has differential
$(Z,Y')\mapsto[Z,Y_{1}]+Y'$ at $(1,Y_{1})$, which is onto because
$\mathfrak{g}_{\RR}=\mathfrak{t}_{1}\oplus[\mathfrak{g}_{\RR},Y_{1}]$ for a
regular semisimple $Y_{1}$. So for neighbourhoods $N_{1}$ of $1$ in
$G(\RR)^{+}$ with $N_{1}u_{1}N_{1}^{-1}\subseteq W$, and $N_{2}$ of $Y_{1}$
in $\mathfrak{t}_{1}$ consisting of regular elements, the image
$\varphi(N_{1}\times N_{2})$ contains a nonempty open subset of
$\mathfrak{g}_{\RR}$. The rational points $\mathfrak{g}(\QQ)$ are dense in
$\mathfrak{g}_{\RR}$, so that open subset contains some
$Y\in\mathfrak{g}(\QQ)$, and $Y=\mathrm{Ad}(g)Y'$ with $g\in N_{1}$ and
$Y'\in N_{2}$. As $Y$ is rational and regular semisimple, the identity
component $T$ of its centraliser in $G$ is a maximal torus defined over
$\QQ$, and $T_{\RR}=gT_{1}g^{-1}$ because the centraliser of $Y'$ has
identity component $T_{1}$. Then $u=gu_{1}g^{-1}$ lies in $W$ and
$u(\mathrm{U}(1))\subseteq T_{\RR}$.
\end{proof}

For an abelian variety the lemma is used with $G$ its Hodge group, the
smallest algebraic $\QQ$-subgroup of $\GL(H^{1}(A,\QQ))$ whose real points
contain the image of the homomorphism $\mathrm{U}(1)\to\GL(H^{1}(A,\RR))$
that defines the Hodge structure. If $\Hg(A)$ lies in a $\QQ$-torus $T$, then
$A$ is of CM type.\footnote{For an elliptic curve $G=\SL_{2}$ and $\cD$ is
the upper half-plane. The compact torus $u_{\tau}(\mathrm{U}(1))$ fixing
$\tau$ lies in the real points of a maximal $\QQ$-torus of $\SL_{2}$ exactly
when $\tau$ is imaginary quadratic, and the torus is then the group of
elements of norm one of $\QQ(\tau)$. So the points of \Cref{lem:cmdense}
are the curves with complex multiplication, and the lemma is the familiar
density of the imaginary quadratic points in the upper half-plane.} Indeed $T$ lies in a maximal $\QQ$-torus $T'$ of
$\GL(H^{1}(A,\QQ))$, which is the group of units of a commutative semisimple
$\QQ$-subalgebra $L\subseteq\End(H^{1}(A,\QQ))$ of dimension $2\dim A$; the
algebra $L$ commutes with $T'\supseteq\Hg(A)$, so it lies in the commutant of
$\Hg(A)$, which is $\End^{0}(A)$.

\begin{proposition}[A family with a CM member through every Hodge
class]\label{prop:cmfamily}
Let $A_{0}$ be a complex abelian variety and $\gamma$ a Hodge class of degree
$2p$ on $A_{0}$. There are a smooth connected complex quasi-projective variety
$S$, a polarised abelian scheme $\pi\colon\cA\to S$, a global section $\xi$ of
$R^{2p}\pi_{*}\QQ$ that is a Hodge class at every point, and points $s_{0}$,
$c\in S$, such that $\cA_{s_{0}}\cong A_{0}$ with $\xi_{s_{0}}=\gamma$ under
this isomorphism and $\cA_{c}$ is of CM type.
\end{proposition}

\begin{proof}
Let $\ell\ge1$ be an integer with $\ell\gamma$ integral. We construct a
section with value $\ell\gamma$ at $s_{0}$ and divide it by $\ell$ at the
end; until then $\gamma$ stands for $\ell\gamma$. Fix a polarisation of
$A_{0}$, put $V=H_{1}(A_{0},\QQ)$, and fix a level $N\ge3$. As in \Cref{setup:shimura}, the
quotient $\cS_{g}$ of the Siegel space $\cH$ of complex structures on
$V_{\RR}$ compatible with the polarisation, by the principal congruence
subgroup of level $N$ of the stabiliser of $H_{1}(A_{0},\ZZ)$, is a smooth
quasi-projective variety carrying a universal family $f\colon\cA_{g}\to\cS_{g}$,
polarised by the fixed polarisation \cite{BB66}; let $a_{0}\in\cS_{g}$ be the
point of $A_{0}$ and $J_{0}\in\cH$ a lift, through which the fibre of
$R^{2p}f_{*}\ZZ$ at every point of the image of $\cH$ is identified with
$H^{2p}(A_{0},\ZZ)=\bigwedge^{2p}H_{1}(A_{0},\ZZ)^{\vee}$ by parallel
transport in the simply connected space $\cH$. Put
$\mathbb{V}=R^{2p}f_{*}\ZZ$, a polarised variation of $\ZZ$-Hodge
structure. By \Cref{thm:cdk} the connected component $\mathcal{C}$ of its
locus of Hodge classes that contains $(a_{0},\gamma)$ is a closed algebraic
subvariety of the total space of $\mathbb{V}$; it has finitely many
irreducible components.

Let $G=\Hg(A_{0})$, which lies in the symplectic group of the polarisation
and fixes $\gamma$ \cite[\S3]{Del82}, let
$u_{0}\colon\mathrm{U}(1)\to G_{\RR}$,
$u_{0}(e^{i\theta})=\cos\theta+J_{0}\sin\theta$, be the homomorphism defining
the Hodge structure, and let
$\cD=\{gJ_{0}g^{-1}:g\in G(\RR)^{+}\}$. For $g\in G(\RR)^{+}$ the complex
structure $J=gJ_{0}g^{-1}$ is compatible with the polarisation and positive
for it, since $g$ preserves the polarisation, so $\cD\subseteq\cH$; and
$\cD$ is the orbit of \Cref{lem:cmdense} when $J$ is identified with
$u_{J}=gu_{0}g^{-1}$. For $J\in\cD$ the group $u_{J}(\mathrm{U}(1))$ lies in
$G(\RR)$, so $\Hg(A_{J})\subseteq G$ fixes $\gamma$, which is therefore a
Hodge class on $A_{J}$. The map sending $J\in\cD$ to the image of $J$ in
$\cS_{g}$ together with $\gamma$ is thus a continuous map from the
connected space $\cD$ to the locus of Hodge classes of $\mathbb{V}$, so it
takes values in $\mathcal{C}$. By \Cref{lem:cmdense} there are points
$J_{k}\in\cD$ converging to $J_{0}$ with $u_{J_{k}}(\mathrm{U}(1))$ in the
real points of a $\QQ$-torus of $G$; then $\Hg(A_{J_{k}})$ lies in that torus
and $A_{J_{k}}$ is of CM type. One of the finitely many irreducible
components of $\mathcal{C}$, say $\mathcal{C}_{0}$, contains the images of
infinitely many $J_{k}$, and, being closed, also their limit $(a_{0},\gamma)$.

Let $\rho\colon S\to\mathcal{C}_{0}$ be a resolution of singularities
\cite{Hir64}: $S$ is smooth, connected and quasi-projective, since
$\mathcal{C}_{0}$ is irreducible and quasi-projective, and $\rho$ is
projective and surjective. Let $\pi\colon\cA\to S$ be the pull-back of
$\cA_{g}$ along $S\to\mathcal{C}_{0}\to\cS_{g}$, a polarised abelian scheme.
A point of $\mathcal{C}_{0}$ is a pair $(a,v)$ with $v\in\mathbb{V}_{a}$ of
Hodge type, and $(a,v)\mapsto v$ is a flat section of the pull-back of
$\mathbb{V}$ to $\mathcal{C}_{0}$, because the total space of a local system
is a covering space of the base; its pull-back to $S$, divided by $\ell$, is
a flat section $\xi$ of $R^{2p}\pi_{*}\QQ$ of Hodge type at every point.
Take $s_{0}$ over $(a_{0},\gamma)$ and $c$ over the image of some $J_{k}$ in
$\mathcal{C}_{0}$.
\end{proof}

\begin{theorem}[(F2) is propagation from CM points]\label{thm:f2propagation}
The following are equivalent.
\begin{enumerate}[label=\textup{(\roman*)},leftmargin=2.6em,itemsep=2pt]
\item \textup{(F2)}.
\item The Hodge conjecture holds for every complex abelian variety.
\item For every polarised abelian scheme $\pi\colon\cA\to S$ over a smooth
  connected complex quasi-projective variety and every global section $\xi$
  of $R^{2p}\pi_{*}\QQ$, if $\xi_{s}$ is algebraic at one point $s\in S$,
  then $\Sigma(\xi)=S$.
\item The same, for the points $s$ at which $\cA_{s}$ is of CM type: if
  $\xi_{c}$ is algebraic at one point $c$ with $\cA_{c}$ of CM type, then
  $\Sigma(\xi)=S$.
\item For every $\pi$ and $\xi$ as in \textup{(iii)} with $\xi_{c}$ algebraic
  at a point $c$ with $\cA_{c}$ of CM type, $\Sigma(\xi)$ has an interior
  point.
\end{enumerate}
\end{theorem}

\begin{proof}
(i) and (ii) are equivalent by \Cref{prop:f2isab}. If (ii) holds and
$\xi_{s}$ is algebraic, then $\xi_{t}$ is a Hodge class at every $t$ by the
first part of the proof of \Cref{prop:lefvhc}, which uses only the theorem of
the fixed part and the semisimplicity of polarisable Hodge structures, and it
is algebraic by (ii); so (ii) implies (iii). (iv) is the case of (iii) at
points of CM type, and
(iv) and (v) are equivalent by \Cref{lem:algstrata}(iii).

Assume (iv). Let $F$ be a CM field and $n\ge2$, and let $\delta_{0}$ be the
trivial class, the discriminant of a hermitian space that is a sum of $n$
hyperbolic planes, each of determinant $-1$. The family of
$(F,n,\delta_{0})$-Weil type over its Shimura variety with a level structure,
as in \Cref{setup:cmfield} and \Cref{setup:shimura}, is a polarised abelian
scheme over a smooth connected quasi-projective base, on which the Weil
classes are global sections (\Cref{prop:cmweil}(i)). By \Cref{thm:main}(i)
they are algebraic at a member isogenous to a product of abelian varieties
with complex multiplication by $F$, which is of CM type, and by (iv) they are
algebraic at every member. Let $A$ be an abelian variety of split Weil type
for $F$ in the sense of \cite[\S2.1]{Mil20}: $F$ acts on it with
$\dim_{F}H^{1}(A,\QQ)=2n$ and $\dim V_{\sigma}^{1,0}=n$ for every embedding
$\sigma$, and for some polarisation the hermitian form of
\Cref{rem:cmmember} has a totally isotropic subspace of dimension $n$, so
that it is a sum of $n$ hyperbolic planes and $A$ is of
$(F,n,\delta_{0})$-Weil type. By \Cref{prop:cmweil}(iv) the Hodge structure
$H^{1}(A,\QQ)$ is isomorphic, compatibly with $F$, to that of a member $A_{s}$
of the family, so there is an isogeny $A\to A_{s}$ commuting with $F$; its
pull-back maps $W_{F}(A_{s})$ onto $W_{F}(A)$ and algebraic classes to
algebraic classes. So every abelian variety of split Weil type has algebraic
Weil classes, for $n\ge2$ by this argument and for $n=1$ by the theorem of
Lefschetz on $(1,1)$ classes. By Andr\'e's theorem
\cite[Theorem~1 and its proof]{Mil20}, every Hodge class on an abelian
variety $B$ of CM type is a sum $\sum f_{\Delta}^{*}(t_{\Delta})$ of
pull-backs along homomorphisms $f_{\Delta}\colon B\to A_{\Delta}$ of Weil
classes $t_{\Delta}$ on abelian varieties $A_{\Delta}$ of split Weil type
for CM fields $F_{\Delta}$; so the Hodge conjecture holds for every abelian
variety of CM type. Now let $A_{0}$ be any abelian variety and $\gamma$ a
Hodge class on it, and take $S$, $\pi$, $\xi$, $s_{0}$, $c$ from
\Cref{prop:cmfamily}. Then $\xi_{c}$ is algebraic, so by (iv)
$\xi_{s_{0}}=\gamma$ is algebraic. This is (ii).
\end{proof}

\Cref{fig:cmseed} draws the theorem in one family.

The implication from \textup{(iii)} to \textup{(ii)} is the argument of
Deligne and Andr\'e as Milne records it \cite[\S5, Theorem~3 and
Remark~2]{Mil20}.\footnote{Deligne's proof that Hodge classes on abelian
varieties are absolutely Hodge uses families of the same kind \cite{Del82}.
By his Principle B \cite[\S2]{Del82} the property of being absolutely Hodge
passes from one member of a family to every member, so a family through a
given Hodge class with a member of CM type reduces the question to that
member. Algebraicity is not known to pass from one member to all, and
\textup{(iii)} is precisely that property.} In that argument the families can be taken to be polarised
abelian schemes and the starting points to be of CM type, which is what
\textup{(iv)} records, and the converse holds. The seeds
are explicit: for the Weil classes they are the base points of
\Cref{thm:cmbasepoint}, and for every other Hodge class they are the CM
members supplied by \Cref{prop:cmfamily}, at which the class is algebraic once
the Weil classes of the split families are. For the Weil classes of the
split family of a CM field, whose base point is of CM type, statement
\textup{(v)} is the condition of \Cref{thm:main}(iv). So the condition of
\Cref{thm:main}(iv), required in every polarised abelian scheme with a
member of CM type at which the class is algebraic, is exactly (F2); it is not a reduction
of (F2) to something smaller, because the interior point it asks for is the
conclusion of \textup{(iii)}.

\begin{figure}[tbp]
\centering
  \includegraphics[width=\textwidth]{fig_cmseed}
\caption{\Cref{thm:f2propagation} in one family. The sheet is the base $S$
of a polarised abelian scheme $\pi\colon\cA\to S$ through a given Hodge
class, as in \Cref{prop:cmfamily}; the small dots are its points of CM type,
dense by \Cref{lem:cmdense}, and the green curves are strata of the algebraic
locus $\Sigma(\xi)$ of the flat class $\xi$ (\Cref{lem:algstrata}). Above
three points the fibres are drawn as tori, and the loop on each stands for
$\xi$, carried from fibre to fibre by parallel transport. At the CM member
$\cA_{c}$ the class is algebraic, and it stays algebraic along the stratum
through $c$. At a very general member $\cA_{s}$ it is algebraic exactly when
$\Sigma(\xi)$ has an interior point, and (F2) is the statement that this
happens in every such family.}
\label{fig:cmseed}
\end{figure}

\subsection{(F3$'$) and the part reached from abelian varieties}
\label{ssec:abpart}

\begin{definition}[The abelian part]\label{def:abpart}
For a smooth complex projective variety $X$ and $p\ge0$ let
$H^{2p}_{\mathrm{ab}}(X)\subseteq H^{2p}(X,\QQ)$ be the sum of the images of
the maps
$\Gamma_{*}\colon H^{k}(B,\QQ)\to H^{2p}(X,\QQ)$,
$y\mapsto\mathrm{pr}_{X*}([\Gamma]\cup\mathrm{pr}_{B}^{*}y)$, over all
complex abelian varieties $B$, the point included, all $k$, and all algebraic
cycles $\Gamma$ on $B\times X$ with rational coefficients and of the
codimension that lands in degree $2p$.
\end{definition}

\begin{proposition}[\textup{(F3$'$)} is a statement about the complement of
the abelian part]\label{prop:abpart}
Let $X$ be a smooth complex projective variety and $p\ge0$.
\begin{enumerate}[label=\textup{(\roman*)},leftmargin=2.6em,itemsep=2pt]
\item $H^{2p}_{\mathrm{ab}}(X)$ is a sub-Hodge structure of
  $H^{2p}(X,\QQ)$, spanned by the images of finitely many of the maps
  $\Gamma_{*}$.
\item $\mathrm{Ab}^{p}(X)=\Hdg^{p}(X)\cap H^{2p}_{\mathrm{ab}}(X)$.
\item Let $T^{2p}_{\mathrm{ab}}(X)$ be a complement of
  $H^{2p}_{\mathrm{ab}}(X)$ in $H^{2p}(X,\QQ)$ as a Hodge structure. Then
  \textup{(F3$'$)} holds for $X$ in degree $2p$ if and only if
  $T^{2p}_{\mathrm{ab}}(X)$ contains no nonzero Hodge class, in particular
  whenever the Hodge group of $T^{2p}_{\mathrm{ab}}(X)$ fixes no nonzero
  vector.
\item $X$ lies in the class $\mathcal{A}$ of \Cref{prop:f3prime}(iii) if and
  only if $T^{2p}_{\mathrm{ab}}(X)=0$ for every $p$ and the odd cohomology is
  reached in the same way.
\end{enumerate}
\end{proposition}

\begin{proof}
(i) An algebraic cycle $\Gamma$ of codimension $e$ on $B\times X$, with
$\dim B=b$, induces a morphism of Hodge structures
$H^{k}(B,\QQ)\to H^{k+2e-2b}(X,\QQ)(e-b)$, so its image is a sub-Hodge
structure; a sum of sub-Hodge structures is one, and a finite-dimensional
space is spanned by finitely many of them.

(ii) Every $\Gamma_{*}\beta$ with $\beta$ a Hodge class lies in both sets, so
$\mathrm{Ab}^{p}(X)$ is contained in the intersection. Conversely take
$(B_{j},\Gamma_{j},k_{j})$, $1\le j\le r$, as in (i), and let
\[
  \Phi=\sum_{j}\Gamma_{j*}\colon
  \bigoplus_{j}H^{k_{j}}(B_{j},\QQ)(t_{j})\longrightarrow
  H^{2p}_{\mathrm{ab}}(X),
\]
with the Tate twists $t_{j}$ that make it a morphism of pure Hodge structures
of weight $2p$; it is surjective. Polarisable Hodge structures form a
semisimple category, so $\Phi$ has a section $\sigma$ that is a morphism of
Hodge structures. For a Hodge class $\gamma\in H^{2p}_{\mathrm{ab}}(X)$ the
class $\sigma(\gamma)$ is of type $(p,p)$, so each component
$\beta_{j}\in H^{k_{j}}(B_{j},\QQ)$ is a Hodge class, $k_{j}$ being even, and
$\gamma=\sum_{j}\Gamma_{j*}\beta_{j}\in\mathrm{Ab}^{p}(X)$.

(iii) A complement exists by semisimplicity, and the Hodge classes of
$H^{2p}(X,\QQ)$ are the sums of those of the two summands; by (ii) the
Hodge classes of the first summand are $\mathrm{Ab}^{p}(X)$. So
$\Hdg^{p}(X)=\mathrm{Ab}^{p}(X)$ exactly when the second has none, and a
Hodge class of $T^{2p}_{\mathrm{ab}}(X)$ is a vector fixed by its Hodge group
\cite[\S3]{Del82}. (iv) is the definition of $\mathcal{A}$ read degree by
degree.
\end{proof}

So \textup{(F3$'$)} asks nothing of the abelian part, where by (ii) and
\Cref{prop:f2isab} the Hodge conjecture reduces to (F2), and it asks of the
complement that it carry no Hodge class. The abelian part contains every
algebraic class, taking $B$ a point, so under the Hodge conjecture the
complement carries none. To prove \textup{(F3$'$)} for $X$ in degree $2p$
without it, one has to show that every Hodge class of $X$ is a sum of
classes $\Gamma_{*}\beta$, that is, to construct algebraic cycles $\Gamma$
on products $B\times X$ that reach it, which is the step
\Cref{rem:f3primeopen} isolates.
\Cref{prop:f3prime}(iii), \Cref{prop:aclosure} and \Cref{thm:simplextype}
are cases in which the complement is zero, and \Cref{prop:k3powers} and
\Cref{prop:k3ntype} cases in which the conjecture itself makes its Hodge
classes algebraic.

\subsection{The first open case of (F3$'$) in families}\label{ssec:rmlocus}

The frontier of \textup{(F3$'$)} in \Cref{rem:f3primefrontier} begins with
$S\times S$ and $S^{[2]}$ for a K3 surface $S$ with real multiplication. These
surfaces come in families over Hermitian symmetric domains, and the argument
of \Cref{thm:main} applies to them with rational Hodge isometries in place of
isogenies.

\begin{setupenv}[A family with real multiplication]\label{setup:rm}
\enlargethispage{2pt}%
Let $E$ be a totally real field of degree $r\ge2$ and $T$ a finite-dimensional
$\QQ$-vector space with an action of $E$ and a nondegenerate symmetric form
$q$ for which $E$ acts by self-adjoint maps, so that $q=\Tr_{E/\QQ}q_{E}$ for
a unique $E$-bilinear form $q_{E}$. Put $m=\dim_{E}T\ge3$, and suppose that
$q_{E}$ has signature $(2,m-2)$ at one real place $\sigma_{1}$ of $E$ and is
negative definite at the others, and that $(T,q)$ embeds isometrically in
$\Lambda_{\QQ}$, where $\Lambda$ is the K3 lattice; write
$N=T^{\perp}\cap\Lambda$. Let $T_{1}=T\otimes_{E,\sigma_{1}}\RR$ and
\[
  \cD_{E}=\{\,[\omega]\in\PP(T_{1}\otimes_{\RR}\CC):\ q(\omega,\omega)=0,\
  q(\omega,\bar\omega)>0\,\}^{+},
\]
one of its two components, a bounded symmetric domain of dimension $m-2$. For
$\omega\in\cD_{E}$ let $S_{\omega}$ be a K3 surface with a marking
$H^{2}(S_{\omega},\ZZ)\cong\Lambda$ carrying $H^{2,0}(S_{\omega})$ to
$\CC\omega$, which exists by the surjectivity of the period map
\cite[Chs.~6--7]{Huy16}; it is projective, since $N$ has signature
$(1,21-rm)$ and so contains a class of positive square, which is of type
$(1,1)$ \cite[Ch.~8]{Huy16}. Each $e\in E$ acts on
$T_{1}\otimes\CC$ by the scalar $\sigma_{1}(e)$ and is self-adjoint, so it
preserves the Hodge decomposition of $T$, and $T\subseteq H^{2}(S_{\omega},\QQ)$
is a sub-Hodge structure containing $T(S_{\omega})_{\QQ}$. Write
$\tilde e_{\omega}\in H^{4}(S_{\omega}\times S_{\omega},\QQ)$ for the class of
the endomorphism $e\oplus0$ of $H^{2}(S_{\omega},\QQ)=T\oplus N_{\QQ}$, a
Hodge class, and let
\[
  \Sigma_{E}=\{\,\omega\in\cD_{E}:\ \tilde e_{\omega}\ \text{is algebraic
  for every}\ e\in E\,\}
\]
be the algebraic locus of the real multiplication; as $e\mapsto\tilde
e_{\omega}$ is $\QQ$-linear, it suffices to ask this for $e$ in a $\QQ$-basis
of $E$. The set $\Sigma_{E}$ does not depend on the choice of the surfaces
$S_{\omega}$. If $(S,\varphi)$ and $(S',\varphi')$ are two marked choices
for $\omega$, then $\psi=\varphi'^{-1}\varphi$ is a rational Hodge isometry
$H^{2}(S,\QQ)\to H^{2}(S',\QQ)$, induced by an algebraic class $Z$ on
$S\times S'$ by Buskin's theorem \cite{Bus19}, and $Z^{t}$ induces the
adjoint $\psi^{-1}$; the composite $Z\circ\tilde e\circ Z^{t}$ acts on
$H^{2}(S',\QQ)$ as $\psi(e\oplus0)\psi^{-1}$, which is $e\oplus0$ for the
marking $\varphi'$, and by zero on $H^{0}$ and $H^{4}$, so it is the class
$\tilde e'$ of $S'$, and $\tilde e$ is algebraic if and only if $\tilde e'$
is. Let $G=\Res_{E/\QQ}\SO(T,q_{E})$, so that
$G(\RR)=\SO(T_{1})\times\prod_{\sigma\ne\sigma_{1}}\SO(T_{\sigma})$; write
$G(\RR)^{+}$ for the identity component of this Lie group, which acts
transitively on $\cD_{E}$ through its first factor, and
$\Gamma_{\QQ}=G(\QQ)\cap G(\RR)^{+}$.
\end{setupenv}

The surfaces with real multiplication of \cite{vG08,vGS25} lie in families
of this kind,\footnote{For a projective K3 surface $S$ the transcendental
part $T(S)_{\QQ}$ is an irreducible Hodge structure whose algebra of Hodge
endomorphisms is a field, totally real or CM \cite{Zar83}; $S$ has real
multiplication when this field is totally real and not $\QQ$. The space
$T(S)_{\QQ}$ is then a vector space over the field, so its degree divides
$\dim T(S)_{\QQ}$, and a very general member of a family of
\Cref{setup:rm} has $T(S_{\omega})_{\QQ}=T$ and real multiplication by
exactly $E$, as the proof of \Cref{thm:rmlocus} shows.} and so do the surfaces $S_{\lambda}$ of \Cref{thm:mumfordks},
for which $E$ is cubic, $m=3$ and $\cD_{E}$ is a disc.

\begin{theorem}[The algebraic locus of a real multiplication]
\label{thm:rmlocus}
Keep \Cref{setup:rm}.
\begin{enumerate}[label=\textup{(\roman*)},leftmargin=2.6em,itemsep=2pt]
\item The points $\omega$ at which $S_{\omega}$ has complex multiplication,
  that is, at which the Hodge group of $T(S_{\omega})_{\QQ}$ is commutative,
  are dense in $\cD_{E}$, and they lie in $\Sigma_{E}$.
\item $\Sigma_{E}$ is stable under $\Gamma_{\QQ}$, and $\Gamma_{\QQ}$ is dense
  in $G(\RR)^{+}$.
\item $\Sigma_{E}$ is dense in $\cD_{E}$, and either $\Sigma_{E}=\cD_{E}$ or
  $\Sigma_{E}$ is meagre.
\item The following are equivalent: $\Sigma_{E}$ is closed; $\Sigma_{E}$ has
  an interior point; $\Sigma_{E}$ is not meagre; $\Sigma_{E}=\cD_{E}$; the
  Hodge conjecture holds for $S_{\omega}\times S_{\omega}$ for every $\omega$
  outside a countable union of proper closed analytic subsets of $\cD_{E}$.
  When they hold, the Hodge conjecture and \textup{(F3$'$)} hold for
  $S_{\omega}\times S_{\omega}$ and $S_{\omega}^{[2]}$ at every $\omega$ at
  which $\End_{\mathrm{Hdg}}(T(S_{\omega})_{\QQ})$ is $E$ or a CM field.
\end{enumerate}
\end{theorem}

\begin{proof}
We first record two facts about the very general point. The Hodge structure
of $S_{\omega}$ on $T$ is given by
$u_{\omega}\colon\mathrm{U}(1)\to\SO(T_{\RR})$, acting by $z^{2}$ on
$\CC\omega$, by $z^{-2}$ on $\CC\bar\omega$ and trivially on their orthogonal
complement; it commutes with $E$, lies in $G(\RR)^{+}$, and
$u_{g\omega}=gu_{\omega}g^{-1}$ for $g\in G(\RR)^{+}$.

(a) \emph{No nonzero $t\in T$ is orthogonal to every $\omega\in\cD_{E}$.} The
lines $\CC\omega$, $\omega\in\cD_{E}$, span $T_{1}\otimes\CC$, since
$\cD_{E}$ is a nonempty open subset of the nondegenerate quadric
$q(\omega,\omega)=0$ in $\PP(T_{1}\otimes\CC)$, so such a $t$
would have zero image in $T_{1}$; but $T\to T_{1}$ is injective, since $T$ is
a free $E$-module and $\sigma_{1}$ is injective.

(b) \emph{An endomorphism $f$ of $T$ that is an endomorphism of Hodge
structures at every $\omega\in\cD_{E}$ is multiplication by an element of
$E$.} Such an $f$ commutes with every $u_{\omega}$, hence with the elements
$X_{\omega}$ of the Lie algebra of $G(\RR)$ that span the Lie algebras of
the groups $u_{\omega}(\mathrm{U}(1))$. They lie in
$\mathfrak{so}(T_{1})\oplus0$, which acts on
$T'=\bigoplus_{\sigma\ne\sigma_{1}}T_{\sigma}$ by zero, and their span is
stable under $G(\RR)^{+}$, since $u_{g\omega}=gu_{\omega}g^{-1}$, so it is an
ideal of $\mathfrak{so}(T_{1})$. That ideal has nonzero projection to every
simple ideal, as the eigenvalues $2i,-2i,0,\dots,0$ of $X_{\omega}$ on
$T_{1}\otimes\CC$ show also for $m=4$, where
$\mathfrak{so}(T_{1})\otimes\CC\cong\mathfrak{sl}_{2}\oplus\mathfrak{sl}_{2}$
and $T_{1}\otimes\CC$ is the tensor product of the two standard
representations; so it is all of $\mathfrak{so}(T_{1})$, and $f$ commutes
with $\mathfrak{so}(T_{1})$. For $m\ge3$ the representation of
$\mathfrak{so}(T_{1})$ on $T_{1}\otimes\CC$ is irreducible and nontrivial, so
$f$ preserves $T_{1}=\mathfrak{so}(T_{1})\cdot T_{\RR}$ and $T'$ and acts on
$T_{1}\otimes\CC=T\otimes_{E,\sigma_{1}}\CC$ by a scalar $\lambda$. An
automorphism $\tau$ of $\CC$, acting on $T\otimes_{\QQ}\CC$ through the
second factor, commutes with $f$ and maps $T\otimes_{E,\sigma_{1}}\CC$ onto
$T\otimes_{E,\tau\sigma_{1}}\CC$, so $f$ acts on the latter by
$\tau(\lambda)$; and every embedding of $E$ is of the form $\tau\sigma_{1}$.
So $f\otimes1$ is multiplication by an element $e$ of $E\otimes_{\QQ}\CC$ on
$T\otimes_{\QQ}\CC=T\otimes_{E}(E\otimes_{\QQ}\CC)$, and $e$ is fixed by
$\mathrm{Aut}(\CC)$ because $f$ is defined over $\QQ$; hence $e\in E$.

For $t\in T$ the set of $\omega$ with $t\perp\omega$ is a closed analytic
subset of $\cD_{E}$, and for $f\in\End_{\QQ}(T)$ so is the set where $f$ is
an endomorphism of Hodge structures, namely where $f(\omega)\in\CC\omega$ and
$f^{*}(\omega)\in\CC\omega$, $f^{*}$ the adjoint. By (a) and (b) these sets
are proper for $t\ne0$ and $f\notin E$. There are countably many of them, and
outside their union $\NS(S_{\omega})_{\QQ}\cap T=0$, so
$T(S_{\omega})_{\QQ}=T$ and $\End_{\mathrm{Hdg}}(T(S_{\omega})_{\QQ})=E$.

(i) The Hodge group of $T$ acts trivially on $T\cap\NS(S_{\omega})_{\QQ}$,
whose orthogonal complement in $T$ is $P=T(S_{\omega})_{\QQ}$, so it is
commutative exactly when the Hodge group $H$ of $P$ is. The group $G$ is
connected and reductive, and $\omega\mapsto u_{\omega}$ identifies $\cD_{E}$
with the orbit of \Cref{lem:cmdense}, since $g\in G(\RR)^{+}$ fixes
$u_{\omega}$ exactly when it preserves the line $\CC\omega$, the eigenspace
of $u_{\omega}(z)$ for $z^{2}$. By \Cref{lem:cmdense} the points at which
$u_{\omega}(\mathrm{U}(1))$ lies in the real points of a maximal $\QQ$-torus of
$G$ are dense; at them the Hodge group of $T$, the smallest algebraic
$\QQ$-subgroup whose real points contain $u_{\omega}(\mathrm{U}(1))$, lies in
that torus, so these points have complex multiplication. Now let $\omega$ be
any point at which $H$ is commutative. The Hodge structure $P$ is
irreducible \cite{Zar83}, so $P$ is an irreducible representation of the
torus $H$, whose rational points are Zariski dense in it. The
$\QQ$-subalgebra $L$ of $\End(P)$ that $H(\QQ)$ generates is therefore
commutative, and $P$ is a simple faithful $L$-module; so $L$ is a field with
$\dim_{L}P=1$, and its commutant $E'=\End_{\mathrm{Hdg}}(P)$ is $L$. By
\cite{Zar83} $E'$ is totally real or CM. If it were totally real, the
embedding $\tau$ through which it acts on the line $H^{2,0}(S_{\omega})$
would be real, and that line would lie in $P\otimes_{E',\tau}\CC$, which is
one-dimensional and stable under complex conjugation, so it would also
contain the different line $H^{0,2}(S_{\omega})$, which is impossible. So
$E'$ is a CM field, the Hodge conjecture
holds for $S_{\omega}\times S_{\omega}$ by \Cref{prop:k3powers}(i), and the
Hodge classes $\tilde e_{\omega}$ are algebraic.

(ii) Let $g\in\Gamma_{\QQ}$ and $\omega\in\Sigma_{E}$. Then $g\omega\in\cD_{E}$,
and $\phi=\mathrm{id}_{N}\oplus g$ is a rational isometry of
$\Lambda_{\QQ}=N_{\QQ}\oplus T$ carrying $\CC\omega$ to $\CC g\omega$, so
through the markings it is a rational Hodge isometry
$H^{2}(S_{\omega},\QQ)\to H^{2}(S_{g\omega},\QQ)$, induced by an algebraic
class $Z\in H^{4}(S_{\omega}\times S_{g\omega},\QQ)$ \cite{Bus19}, whose
transpose induces $\phi^{-1}$. As $g$ is $E$-linear,
$\phi\circ(e\oplus0)\circ\phi^{-1}=e\oplus0$, so, as in \Cref{setup:rm},
$\tilde e_{g\omega}=Z\circ\tilde e_{\omega}\circ Z^{t}$ is algebraic.

For the density of $\Gamma_{\QQ}$, let $Y$ be an $E\otimes\RR$-linear map of
$T\otimes\RR$ that is skew-adjoint for $q$ and has $1-Y$ invertible. Its
Cayley transform $c(Y)=(1+Y)(1-Y)^{-1}$ is $E\otimes\RR$-linear, has
$c(Y)^{*}c(Y)=1$ because $1-Y$ is the adjoint of $1+Y$, and has determinant
one for the same reason, so $c(Y)\in G(\RR)$; and $c$ is a homeomorphism onto
the open neighbourhood $\{g\in G(\RR):1+g\ \text{invertible}\}$ of $1$, with
inverse $g\mapsto(g-1)(g+1)^{-1}$. The $E$-linear skew-adjoint maps of $T$
are dense among these $Y$, and $c$ maps them into $G(\QQ)$. So $\Gamma_{\QQ}$
is dense in a connected open neighbourhood of $1$, its closure is an open
subgroup of $G(\RR)^{+}$, and it is all of $G(\RR)^{+}$.

(iii) and (iv). Density follows from (i). If $\Sigma_{E}$ contains a nonempty
open set $U_{0}$, then $\Sigma_{E}=\cD_{E}$: for $\omega\in\cD_{E}$ choose
$\omega'\in U_{0}$ and $g_{1}\in G(\RR)^{+}$ with $g_{1}\omega'=\omega$; the
set of $g\in G(\RR)^{+}$ with $g^{-1}\omega\in U_{0}$ is open and contains
$g_{1}$, so by (ii) it contains some $g\in\Gamma_{\QQ}$, and
$\omega=g(g^{-1}\omega)\in\Sigma_{E}$ by (ii).

It remains to show that $\Sigma_{E}$ has an interior point as soon as it is
not meagre. Fix $h\in N$ with $h^{2}>0$ and $h\cdot\delta\ne0$ for every
$\delta\in N$ with $\delta^{2}=-2$, which exists because the hyperplanes
$\delta^{\perp}$ are locally finite in the positive cone of $N_{\RR}$
\cite[Ch.~8]{Huy16}. Let $\cD^{\circ}$ be the set of $\omega\in\cD_{E}$ that
are orthogonal to no $\delta\in\Lambda$ with $\delta^{2}=-2$ and
$h\cdot\delta=0$. Writing such a $\delta$ as $n+t$ with $n\in N_{\QQ}$ and
$t\in T$, we have $t\ne0$, because $N_{\QQ}\cap\Lambda=N$, and
$\delta\perp\omega$ exactly when $t\perp\omega$; so by (a) the complement of
$\cD^{\circ}$ lies in a countable union of proper closed analytic subsets and
is meagre.

Let $\omega\in\cD^{\circ}$. No class of $\NS(S_{\omega})$ of square $-2$ is
orthogonal to $h$, so after composing the marking with an isometry $\pm w$,
$w$ a product of reflections in such classes, which fixes $\omega$ and so
changes only the choice of $S_{\omega}$, the class $h$ is that of an ample
line bundle $L$ on $S_{\omega}$ \cite[Ch.~8]{Huy16}. The deformations of the
pair $(S_{\omega},L)$ form a smooth family over a base that the local period
map identifies with a neighbourhood of $\omega$ in the domain of periods
orthogonal to $h$ \cite[Ch.~6]{Huy16}, and that domain contains $\cD_{E}$,
since $h\in N$. Restricting to $\cD_{E}$, we obtain an open neighbourhood $U$
of $\omega$ in $\cD_{E}$ and a smooth family $\cX_{U}\to U$ of K3 surfaces with
markings, whose fibre over $u$ is a choice of $S_{u}$, carrying a line bundle
that restricts to $L$ on $S_{\omega}$. On $S_{\omega}$ the bundle
$L^{\otimes3}$ is very ample and has no higher cohomology
\cite[Ch.~2]{Huy16}; by semicontinuity and base change, after shrinking $U$,
its direct image is free and its sections embed $\cX_{U}$ in
$\PP^{M}\times U$ over $U$. The classes $\tilde e$ are flat sections of the
fourth cohomology of $\cX_{U}\times_{U}\cX_{U}$, which is projective over $U$.
As in the proof of \Cref{lem:algstrata}(i), with the Hilbert schemes of
$\PP^{M}\times\PP^{M}$ \cite{Gro61}, which are projective, for each $e$ in a
$\QQ$-basis of $E$ the set of $u\in U$ at which
$\tilde e_{u}=\sum_{i}q_{i}[Z_{i}]$ with rational $q_{i}$ and $Z_{i}$ of
Hilbert polynomial $P_{i}$ is the image, under a proper map, of a union of
connected components of a fibre product of relative Hilbert spaces over $U$,
and so it is closed in $U$. Taking the union over the countably many
$(\underline{P},\underline{q})$ and the intersection over the basis,
$\Sigma_{E}\cap U$ is a countable union of subsets of $U$ that are closed in
$U$.

The sets $U$ constructed for the points of $\cD^{\circ}$ cover it, and
countably many of them, $U_{i}$, suffice, since $\cD_{E}$ has a countable
basis. If $\Sigma_{E}$ is not meagre, neither is
$\Sigma_{E}\cap\cD^{\circ}$, so some $\Sigma_{E}\cap U_{i}$ is not meagre, and
by the Baire category theorem in $U_{i}$ one of the countably many closed
subsets of $U_{i}$ that make it up has an interior point. This proves (iii),
and in (iv) the equivalence of the first four conditions, since a dense
closed set is everything. The fifth condition implies the third, because at
each point where the Hodge conjecture holds for $S_{\omega}\times S_{\omega}$
the Hodge classes $\tilde e_{\omega}$ are algebraic, and the complement of a
countable union of proper closed analytic subsets is not meagre.

If $\Sigma_{E}=\cD_{E}$ and $\End_{\mathrm{Hdg}}(T(S_{\omega})_{\QQ})=E$, then
$E$ preserves $T(S_{\omega})_{\QQ}$, the smallest sub-Hodge structure of $T$
whose complexification contains $H^{2,0}(S_{\omega})$, since each $e\ne0$
maps it onto a sub-Hodge structure of the same dimension that contains
$H^{2,0}(S_{\omega})$; and each $e\in E$ acts
on it through the algebraic class $\pi_{T}\circ\tilde e_{\omega}\circ\pi_{T}$,
where $\pi_{T}$ is the algebraic projector onto $T(S_{\omega})_{\QQ}$ of the
proof of \Cref{prop:k3powers}; so \Cref{prop:k3powers}(iii) gives the Hodge
conjecture for $S_{\omega}\times S_{\omega}$ and $S_{\omega}^{[2]}$, and with
it \textup{(F3$'$)}. If the field is CM, \Cref{prop:k3powers}(i) gives the
same. Outside the countable union found at the beginning of the proof the
field is $E$, so the fourth condition implies the fifth.
\end{proof}

On the families of van Geemen and Sch\"utt with real multiplication by
$\QQ(\zeta_{7}+\zeta_{7}^{-1})$ and $\QQ(\zeta_{9}+\zeta_{9}^{-1})$ the
real multiplication is induced by explicit cycles \cite[Theorem~1.2 and
Section~4.8]{vGS25}, and on the double planes branched along six lines that
have real multiplication, necessarily by a quadratic field here, by
Schlickewei's cycles \cite{Sch10,vG08}. Whenever the members on which such
cycles are known contain a nonempty open subset of $\cD_{E}$, (iv) gives
$\Sigma_{E}=\cD_{E}$. For the family of the surfaces $S_{\lambda}$,
\Cref{thm:rmlocus} gives the real multiplication at every CM point and on a
dense set stable under $\Gamma_{\QQ}$, and \Cref{thm:mumfordks}(ii) gives it
at $S_{\lambda}$ if the Kuga--Satake class of $S_{\lambda}$ is algebraic;
whether $\Sigma_{E}$ is closed there is open. Since $\Sigma_{E}$ is stable
under the group $\Gamma_{\QQ}$, every $\Gamma_{\QQ}$-orbit lies either in
$\Sigma_{E}$ or in its complement: the Hodge isometries of (ii) move the real
multiplication along an orbit and never bring it to an orbit where it is not
yet algebraic, and by \Cref{rem:mumfordks} the composites of Hodge
similarities that return to $T(S_{\lambda})$ span only the maximal
multiquadratic subfield of $E$, which is $\QQ$ for a cubic field.

\subsection{What the constructions reach}\label{ssec:reach}

\begin{theorem}[Where (F2) and (F3$'$) stand]\label{thm:twostand}
\leavevmode
\begin{enumerate}[label=\textup{(\roman*)},leftmargin=2.6em,itemsep=3pt]
\item \emph{What carries over.} In every polarised abelian scheme of
  \Cref{thm:f2propagation} the algebraic locus of a flat class is a
  countable union of closed algebraic strata and is either everything or
  meagre (\Cref{lem:algstrata}); in the Weil families it is moreover dense
  and contains an explicit CM point (\Cref{thm:main}); in every family of
  \Cref{setup:rm} the algebraic locus of the real multiplication is dense,
  contains every CM point, is stable under $\Gamma_{\QQ}$ and is either
  everything or meagre (\Cref{thm:rmlocus}); and every Hodge class of every
  abelian variety lies in a polarised abelian scheme with a CM member
  (\Cref{prop:cmfamily}).
\item \emph{(F2).} It holds if and only if the algebraic locus of every flat
  Hodge class algebraic at a CM member of a polarised abelian scheme has an
  interior point (\Cref{thm:f2propagation}). In the Weil families the
  constructions of the paper that supply an interior point are: an object
  with injective semiregularity map, or satisfying the flatness condition,
  at one point (\Cref{thm:semiregclosure}, \Cref{thm:cmpropagation},
  \Cref{prop:p2flat}); a bound on the Hilbert data of representatives along
  the Hecke orbit of a base point (\Cref{thm:degreebound}), which is
  equivalent to its conclusion (\Cref{thm:p1equivalent}); and the Lefschetz
  standard conjecture for the total space over a curve through a base point,
  which implies it and, when the monodromy is Zariski dense, follows from it
  (\Cref{cor:weilequiv}). For the squares of Mumford fourfolds the last two
  take the forms of \Cref{thm:modpdensity} and \Cref{cor:mumfordequiv}.
  Markman's sheaves are objects of the first kind, and with them he proves
  the algebraicity of the Weil classes on every abelian fourfold of Weil
  type and on the split sixfolds of every imaginary quadratic field
  \cite{Mar25a,Mar25b}; with \cite{MZ99} this gives the Hodge conjecture for
  abelian varieties of dimension at most five. The open cases named in
  \Cref{sec:scope} include the non-split Weil sixfolds of imaginary
  quadratic fields, the Weil eightfolds of quartic CM fields, where an
  object of the first kind needs
  $\dim\Ext^{2}(E,E)\ge88$ (\Cref{cor:quarticleast}), and the squares of
  Mumford fourfolds, where the classes are algebraic at every CM point
  (\Cref{thm:mumfordcm}) and on every member once one Kuga--Satake class is
  algebraic (\Cref{thm:mumfordks}), or once one complex at one CM point has
  $\Ext^{2}$ of dimension $119$ (\Cref{thm:mumfordnumerical}). With the
  theorems of the preprints of \Cref{sec:preprints} as hypotheses, the
  non-split sixfolds and the squares of Mumford fourfolds are settled
  (\Cref{thm:weilsix}, \Cref{cor:mumfordall}), and the first open cases are
  the eightfolds of nontrivial discriminant and the Weil eightfolds of
  quartic CM fields (\Cref{rem:remains}).
\item \emph{(F3$'$).} It holds if and only if, for every variety and every
  degree, the complement of the abelian part carries no Hodge class
  (\Cref{prop:abpart}). The complement is zero on the class $\mathcal{A}$,
  which contains the Fermat and Delsarte hypersurfaces of every dimension and
  the cyclic covers branched along them (\Cref{prop:aclosure},
  \Cref{cor:delsarteall}), and it carries no Hodge class on the varieties of
  \Cref{prop:f3primesmall}, \Cref{prop:f3primefibration} and
  \Cref{prop:f3primedominant}, where \textup{(F3$'$)} is proved, and of
  \Cref{prop:k3powers} and \Cref{prop:k3ntype}, where the conjecture itself
  is.
  On the first open varieties, $S\times S$ and $S^{[2]}$ for a K3 surface
  with real multiplication, the algebraic locus of the real multiplication
  is dense, contains every CM point and is either the whole family or meagre
  (\Cref{thm:rmlocus}); the real multiplication is algebraic on the members
  of the families of \cite{vGS25} and on the quadratic double planes of
  \cite{Sch10}, and whether the locus is closed is open for the family of
  $S_{\lambda}$; with the preprints of \Cref{sec:preprints} as hypotheses
  the locus is the whole family in every case (\Cref{cor:k3all}). Beyond the
  varieties of abelian type, a Hodge
  class off the abelian part is removed only by a new algebraic cycle on a
  product $B\times X$ (\Cref{rem:f3primeopen}).
\end{enumerate}
Neither (F2) nor \textup{(F3$'$)} is proved here, and by \Cref{thm:final}
together they are the Hodge conjecture.
\end{theorem}

\begin{proof}
Each statement is the result cited with it. For (ii), Markman's theorems are
recalled in the introduction and their reach in \Cref{sec:scope}, where the
non-split sixfolds and the fields of degree at least four are among the open
entries, and \Cref{thm:quarticrank} and
\Cref{cor:quarticleast} give the number $88$. For (iii), the known cases of
the real multiplication are those recalled in \Cref{rem:mumfordks}.
\end{proof}
